\begin{corollary}\label{cor:regular-pair}
  If $R \subseteq \mathbb{Z}^2$ is simple, then $R$ does not admit any pair-free domino tilings.
\end{corollary}

\subsection{Divergence of a domino tiling}

In general, for a subset $R \subseteq \mathbb{Z}^d$ and a tiling $T$ of $R$, we may  define $\delta_T \colon R \to T$ to be the map taking a point to the unique domino containing it, and define the \emph{divergence} of $T$ by
\begin{equation}
  (\operatorname{div} T)(Q) = \frac{1}{2^{d}}\sum_{q \in Q} \iota(\delta_T(q); Q).
\end{equation}
The quantity's name is justified as follows. First, define $\sigma_T \colon R \to R$ to be the map taking a point to the other point in its containing domino. Then, if one realizes $R$ as the union of closed cubes
\begin{equation}
\Omega_R = \bigcup_{q \in R} \overline{B_{1/2}(q)},
\end{equation}
and defines the vector field
\begin{equation}
X_T(p) = 
\begin{cases}
  \sigma_T(q) - q &\text{if } p \in B_{1/2}(q) \text{ for some } q \in R \\
  0 &\text{otherwise}
\end{cases},
\end{equation}
then the divergence at a lattice cube $Q \cong \{0,1\}^d$ centered at $p \in \frac{1}{2}\mathbb{Z}^d$ is given by
\begin{equation}
(\operatorname{div} T)(Q) = \lim_{s \to 1^{-}} \, \frac{1}{\operatorname{vol}_{d-1} (\partial B_{2s}(p))} \int_{\partial B_{2s}(p)} X_T \cdot \mathbf{n}
\end{equation}
This may be compared with divergence of a smooth vector field, in which case the limit is to be taken as $s \to 0$.